\section{Discussion}\label{sec:disc}

This study demonstrates that using several obligate hosts does not by itself stabilize a parasite against host fluctuations, whereas spending part of the life cycle in a stage sheltered from the fluctuating transfer does reduce the long-run growth loss under stated conditions. \citet{dobson2020} proposed that complex life cycles carry an intrinsic stability for two reasons, first because geometric averaging across $n$ hosts takes an $n$th root of the product of transmission terms, and second because long-lived stages store the parasite across another host's booms and busts. When growth is measured per unit time rather than per step of the next-generation matrix, the $n$th-root effect disappears in both deterministic regulation and stochastic persistence (Sections~\ref{sec:model} and~\ref{sec:ring}), consistent with the general equivalence of threshold criteria across life-cycle representations \citep{rueffler2013}. Adding obligate hosts in sequence without lengthening the generation time leaves two stages in one host and two hosts in series mathematically identical whenever their joint law of residence time and survival is the same (Section~\ref{sec:mech}). In contrast, the storage intuition holds as a statement about residence times. In a primitive ring with memoryless stays and scarcity that wastes failed contacts, the second-order loss in $\ln R_0$ per generation scales with log-host variance inflated by positive host memory and divided by the mean generation time $\mathcal{T}$, so a long sheltered stay drives that coefficient toward zero (Section~\ref{sec:ring}), in line with demographic buffering in structured free-living populations \citep{tuljapurkar1982a,tuljapurkar1990,morris2008,jonsson2021}. In the two-stage gamma-host model solved by \citet{roerdink1988}, we proved that this reduction holds strictly at every stay length and every noise amplitude, with the small-noise coefficient becoming asymptotically exact for long stays (Section~\ref{sec:exact}). Whether multiple hosts add buffering therefore depends on how they enter the life cycle rather than on host count alone. Obligate hosts in series do not protect by their number, whereas at fixed mean total transfer, equal shares of alternative hosts with matched log variance and memory used in parallel at one step divide that step's second-order loss by their number when their fluctuations are independent, give a weaker reduction under partial synchrony, provide no reduction in perfect synchrony, and can raise the loss when a newcomer is much noisier (Proposition~\ref{prop:ring-alt}, Section~\ref{sec:example}).

\subsection{Biological limits of stage-based buffering}\label{subsec:disc-limits}

Where that protection stops is set by life-history biology rather than mathematical technicalities (Sections~\ref{sec:mech} and~\ref{sec:num}; assumption ladder A1--A13 in Online Resource~1). First, positive host memory slows the benefit of storage and can reverse it at short durations. In a primitive two-stage geometric ring with exposed stay $a$ and steady stay $b$, adding $b \ge 1$ raises the second-order loss above the one-stage value when $a > 1$ and $b < a(2a - 1)\theta/(a + \theta)$ for memory scale $\theta = \phi/(1 - \phi)$, while lengthening $b$ can raise the coefficient under strong memory when $a > (3 + \sqrt{17})/4$ (and in a three-stage $(1, b, c_3)$ ring when $c_3 > 2$, matching nonmonotonicity in 15 of 24 finite-noise three-stage grids; Section~\ref{sec:exact}). Second, the shape of the residence-time distribution governs how lineages spread across host states \citep{tuljapurkar1990,lloyd2001}. At fixed mean generation time under positive geometric memory ($0 < \phi < 1$), greater dispersion in circuit times lowers the second-order loss through $\mathbb{E}[\phi^L]$ and can produce a net gain under heavy-tailed dormancy or relapse \citep{white2016}, whereas regular developmental clocks erode buffering until an all-fixed cycle under $n$-step ergodicity suffers the entire logarithmic-gap penalty $D = \ell$ (Section~\ref{sec:mech}). Third, survival during host scarcity and storage bounds what a long stage can save. When scarcity delays contact without waiting mortality ($d = 0$), reproductive value is conserved and $D = 0$ at every admissible amplitude, whereas waiting deaths ($0 < d \le 1$) leave a positive second-order floor $\sigma^2 d(1 - d)q_j$ even as another sheltered stay grows without bound (Proposition~\ref{prop:mech-contact}). Mortality within the sheltered stage itself taxes baseline reproduction, and in finite whole-number populations, parasites grouped in packets ($K = 100$) that share their host's death ($a = 1$) still experienced $0.1960 \pm 0.0126$ extinction at stay $160$ across mean-$200$ outages despite positive long-run growth (Section~\ref{sec:num}).

Shared or cyclic host forcing and complete host absences alter even the sign of the steady-stage effect. When a single fluctuating vector or host mediates two transfers in the loop \citep{watts2021}, or when two obligate hosts covary as predator and prey, cross-lag terms enter the second-order loss (Proposition~\ref{prop:mech-linked}). In discrete time, one host moving two unequal transfers together turns fluctuations into a second-order gain above a memory threshold (such as $\phi > 0.253$ for stays $(1, 10)$), and for linked predator--prey hosts the sequence in which the parasite uses them around the cycle changes the outcome, shifting $D$ from $+0.201$ to $-0.059 \pm 0.002$ at $\sigma = 1$ for stays $(5, 50, 5)$ (Section~\ref{sec:mech}). Similarly, when host numbers alternate or cycle in step with the parasite's generation time \citep{tuljapurkar1985,botsford2014}, resonance can make $D$ negative (Proposition~\ref{prop:ring-spectrum}), so lengthening a sheltered stage erodes a reproductive gain rather than softening a loss ($D_{\mathrm{alt}}'(s) > 0$ for $0 \le s < 1$ under period-two alternation). A parallel reversal arises when an obligate host drops to zero (assumption A13), where logarithmic moments diverge ($\ell = \infty$). Whereas persistent random outages keep $D \ge 0$ under temporal independence and leave a generally positive long-stay bound limit $-\ln[(1 - q_0)G_+]$ whenever positive host states vary, period-two boom-and-zero cycles $h = (2, 0)$ yield $D < 0$ with $D_s > 0$ at every stay, and at long stays the fourth-moment raw-covariance sum $A = \operatorname{Var}(h)/2 + \sum_{\ell \ge 1} p^{\ell - 1}\operatorname{Cov}(h_t, h_{t-\ell})$ sets the sign of $D$ and of fixed-factor stay comparisons (Section~\ref{sec:mech}). Once an established parasite depresses its hosts (Section~\ref{sec:dyn}), stability shifts again to endemic equilibrium damping \citep{anderson1978,dobson1992}, where external larval delays can destabilize weakly self-regulated hosts and none of 810 baseline trophic equilibria without prey self-limitation was stable. These boundaries are biological mechanisms of real transmission cycles, not technical caveats around a universal rule.

\subsection{Scope and asymmetry of the evidence}\label{subsec:disc-evidence}

A sharp asymmetry separates what this study establishes mathematically from what existing empirical records can decide. On the theoretical side, the analytical propositions and theorems in Sections~\ref{sec:model}--\ref{sec:dyn} rest on written proofs audited by independent second-party derivations and hostile referee checks. Across the 82 theorem, proposition, lemma, and corollary environments in the main text and appendices (25 in Sections~\ref{sec:model}--\ref{sec:dyn} and 57 in the appendices), Lean~4 checks the complete statements of nine (two in the main text), meaning every clause from the stated model, with named textbook limit theorems such as the Furstenberg--Kesten limit as explicit hypotheses. For 57 others (16 in the main text) it checks specified algebraic or analytic components, and for the remaining 16 (seven in the main text) it checks components that still carry written proof obligations. Random-matrix Lyapunov limits, stationary Markov measures, and the remaining analytic and asymptotic steps rely on the written proofs. On the simulation side, an independent implementation reproduced all 190 of 190 validation-gate comparisons across five gates and all 18 pilot cells of the demographic world, and driving that world with 120 open population series and 71 \textit{Tribolium} records traced all 271 primary departures from baseline analytical formulas without free parameters, assigning 220 to existing finite-amplitude, higher-moment, or spectral-resonance results, 51 to the absent-host extension A13, and 0 to unresolved causes (Section~\ref{sec:num}). In contrast, the data did not decide whether stage buffering stabilizes real parasites. Across seven pre-registered case studies spanning museum fish parasites, 1,781 host series, bacteria--phage vessels, microparasite microcosms and lakes, tick-borne \textit{Borrelia} surveys, cyclic voles, and 383 strict stage-duration cohorts (Table~\ref{num-tab-casestudies}), open records confirm that positive host memory, large fluctuation amplitudes, oscillatory rodent cycles, and non-exponential stage durations are realistic features of natural systems. However, none of the seven case studies drew on data that vary exogenous host fluctuations at fixed arithmetic-mean host abundance while measuring stage-duration laws and multi-generation persistence, leaving every direct empirical test undecided or confounded by mean-host shifts, thermal lethality, resistance evolution, and endogenous host crashes.

\subsection{Predictions for control, climate, and parasite decline}\label{subsec:disc-ecology}

Because the empirical studies remain inconclusive, ecological consequences of the theory must be stated as testable predictions rather than established field findings. First, for infectious disease control, we predict that pulsed vector suppression or intermittent transmission interruption at a fixed time-averaged intensity will penalize short, synchronous life cycles far more strongly than cycles anchored by long, dispersed sheltered stages \citep{smith2007,colley2014,anderson2015,gurarie2018}. In the models, multi-year adult helminths, overdispersed asymptomatic blood-stage infections \citep{andrade2020}, and relapsing liver stages \citep{white2016} absorb vector troughs, and the outage retention law in Eq.~\eqref{num-eq-gap} predicts that local elimination under buffered stages requires either maintaining zero-transmission gaps across multiple mean reservoir lifetimes or coupling vector pulses with treatment of the sheltered stage itself. Second, under climate change, we predict that complex life cycles face two interacting channels of risk \citep{morris2008,molnar2013,kutz2014,aleuy2020}. Greater rainfall and temperature variability raises host fluctuation variance $\sigma^2$ and outage frequency $q_0$, while warming accelerates ectothermic development (temperature correlated with mean stage duration at $\rho = -0.943$ across 479 cohorts in Table~\ref{num-tab-casestudies}, with $R^2 = 0.0032$ on individual CV), which shortens generation time $\mathcal{T}$ and is predicted to strip away temporal buffering unless endotherm stages supply a long, thermally insulated stay or opposing thermal slopes of mortality and development offset the loss (Proposition~\ref{prop:mech-thermal}). Third, for biodiversity loss and long-term parasite declines \citep{lafferty2009,lafferty2012,lagrue2015,wood2023}, we predict that obligate host count is a poor proxy for fluctuation risk once permanent host extirpation is separated from boom-and-bust dynamics. Instead of counting host species, field monitoring should measure sheltered and exposed stage-duration distributions, waiting-death fractions $d$ during host scarcity, directed lag covariances between predator and prey hosts, and the frequency of local host absences.

\subsection{Experimental test and open problems}\label{subsec:disc-open}

Resolving the empirical gap requires a manipulative laboratory experiment that crosses fixed-mean host supply (constant versus fluctuating host inflow at identical time-averaged host density, run in complete mirrored schedule pairs) with stage duration outside the fluctuating transfer (Section~\ref{sec:num}). In our corrected demographic simulations on matched reference-generation horizons, detecting the buffering interaction at $80\%$ power required 51--61 independent replicate populations per cell across the three valid primary horizons, or ca.\ 244 independent cultures ($61 \times 4$) for a four-cell core comparison. Candidate host--parasite systems with short generation times and free transmission stages include \textit{Paramecium caudatum} with \textit{Holospora undulata} \citep{duncan2013}, bacteria with lytic phage \citep{harrison2013}, and \textit{Daphnia} with \textit{Ordospora} \citep{krichel2023}, while parasite-free \textit{Tribolium castaneum} microcosms show that fixed-mean resource schedules can be maintained experimentally \citep{desharnais2023}. However, none of these systems is ready for an uncalibrated test because dilution or harvesting can remove free stages alongside hosts, rapid phage--bacteria runs evolve resistance, and parameter-free predictions require prior assays of encounter curvature, waiting mortality $d$, both stage-duration distributions, starting stage composition, and observation error.
